\begin{abhakseite}
\ifmitzone
\abhakgruppe{Kennst du schon}
\abhak{1}{Ich kann negative Zahlen addieren und subtrahieren.}
\abhak{2}{Ich kann Zahlen mit Vorzeichen malnehmen.}
\abhak{3}{Ich kann Punkt vor Strich rechnen.}
\abhak{4}{Ich kann den Wert eines Terms berechnen.}
\abhak{5}{Ich kann gleichartige Glieder zusammenfassen.}
\abhak{6}{Ich kann Terme malnehmen.}
\abhak{7}{Ich finde den Fehler beim Zusammenfassen.}
\abhak{8}{Ich kann Terme mit $x$ und $x^2$ richtig zusammenfassen.}
\abhak{9}{Ich kann Quadratzahlen und Quadratwurzeln angeben.}
\fi
\abhakgruppe{Einheit 1 · Klammern auflösen}
\abhak{10}{Ich erkenne, welches Vorzeichen zu welchem Glied gehört.}
\abhak{11}{Ich erkenne, was vor einer Klammer steht.}
\abhak{12}{Ich kann Klammern auflösen (neu in diesem Jahr).}
\abhak{13}{Ich kann Klammern auflösen und den Term zusammenfassen.}
\abhak{14}{Ich finde den Fehler beim Auflösen einer Klammer.}
\abhak{15}{Ich kann begründen, ob zwei Terme gleichwertig sind.}
\abhak{16}{Ich kann eine Sachaufgabe mit Klammern lösen.}
\abhakgruppe{Einheit 2 · Ausklammern}
\abhak{17}{Ich kann einen gemeinsamen Faktor ausklammern.}
\abhak{18}{Ich finde den Fehler beim Ausklammern.}
\abhak{19}{Ich kann entscheiden, ob alle Glieder einen gemeinsamen Faktor haben.}
\abhak{20}{Ich kann eine Sachaufgabe durch Ausklammern lösen.}
\abhakgruppe{Einheit 3 · Summe mal Summe}
\abhak{21}{Ich erkenne, was mit was malgenommen wird.}
\abhakgruppe{\normalfont\itshape Terme mit zwei Variablen}
\abhak{22}{Ich kann den Wert eines Terms mit zwei Variablen berechnen.}
\abhak{23}{Ich kann Terme mit zwei Variablen zusammenfassen (neu in diesem Jahr).}
\abhak{24}{Ich kann eine Klammer mit zwei Variablen auflösen.}
\abhakgruppe{\normalfont\itshape Klammer mal Klammer}
\abhak{25}{Ich kann zwei Klammern miteinander malnehmen.}
\abhak{26}{Ich finde den Fehler beim Malnehmen zweier Klammern.}
\abhak{27}{Ich kann am Rechteck begründen, warum vier Produkte entstehen.}
\abhak{28}{Ich kann eine Sachaufgabe mit zwei Klammern lösen.}
\abhakgruppe{Einheit 4 · Binomische Formeln}
\abhak{29}{Ich erkenne, ob eine Klammer mit sich selbst malgenommen wird.}
\abhak{30}{Ich erkenne, welche binomische Formel passt.}
\abhak{31}{Ich erkenne, was $a$ und was $b$ ist.}
\abhakgruppe{\normalfont\itshape Ausmultiplizieren mit den Formeln}
\abhak{32}{Ich kann mit den binomischen Formeln ausmultiplizieren.}
\abhak{33}{Ich kann mit den binomischen Formeln ausmultiplizieren – weiter.}
\abhakgruppe{\normalfont\itshape Zwei Terme als gleich nachweisen}
\abhak{34}{Ich kann nachweisen, dass zwei Terme gleich sind.}
\abhakgruppe{\normalfont\itshape Kopfrechnen mit den Formeln}
\abhak{35}{Ich kann mit den binomischen Formeln im Kopf rechnen.}
\abhak{36}{Ich finde den Fehler beim Ausmultiplizieren mit einer binomischen Formel.}
\abhak{37}{Ich kann entscheiden, welcher Term gleichwertig ist.}
\abhak{38}{Ich kann begründen, warum beim Quadrat einer Summe ein Mittelglied entsteht.}
\abhak{39}{Ich kann eine Sachaufgabe mit einer binomischen Formel lösen.}
\abhakgruppe{Einheit 5 · Faktorisieren}
\abhak{40}{Ich erkenne, ob ein Term ein Quadrat ist.}
\abhak{41}{Ich erkenne die beiden Quadrate und das Mittelglied.}
\abhakgruppe{\normalfont\itshape Faktorisieren mit den Formeln}
\abhak{42}{Ich kann einen Term mit einer binomischen Formel faktorisieren.}
\abhakgruppe{\normalfont\itshape Gleichungen durch Faktorisieren lösen}
\abhak{43}{Ich kann eine Gleichung lösen, indem ich faktorisiere.}
\abhakgruppe{\normalfont\itshape Quadratische Ergänzung}
\abhak{44}{Ich kann einen Term als Quadrat einer Klammer plus eine Zahl schreiben.}
\abhak{45}{Ich finde den Fehler beim Faktorisieren.}
\abhak{46}{Ich kann entscheiden, ob sich ein Term mit einer binomischen Formel faktorisieren lässt.}
\abhak{47}{Ich kann begründen, wann ein Term keine binomische Formel ist.}
\abhak{48}{Ich kann eine Sachaufgabe durch Faktorisieren lösen.}
\end{abhakseite}
